\chapter{网络流}

\section{最大流问题}

\section{费用流}

% \begin{summary}
% \begin{itemize}
%     \item 分治信息的概念与设计原则：首先将复杂的信息分解为相对简洁的信息，再从简洁的信息开始逐步添加合并信息所需的辅助信息，直到信息能够合并。
%     \item 静态区间数据结构的设计方法：倍增（ST 表）与分治（二区间合并、线段树、多叉分治、一般的分治结构）
%     \item 线段树的结构与区间分解过程
% \end{itemize}
% \end{summary}


% \begin{table}[h]
% \centering
% \begin{tabular}{@{}llll@{}}
% \toprule
% 数据结构 & 预处理 & 查询 & 适用范围 \\ \midrule
% 前缀和 & $O(n)$ & $O(1)$ & 具有逆元 (群) \\
% ST 表 & $O(n \log n)$ & $O(1)$ & 可重复贡献 (半格) \\
% 线段树 & $O(n)$ & $O(\log n)$ & 一般半群 \\
% Disjoint ST & $O(n \log n)$ & $O(1)$ & 一般半群 \\
% \bottomrule
% \end{tabular}
% \end{table}